\begin{exercisechunk}
\item \textbf{Continuous EM and reference measures.}
\label[exercise]{ex:l5-continuous-em-reference}
\exerciseprofile{3}{3}{3}{\exproof\quad\excalculation\quad\exinterpretation}{%
Prove continuous-EM monotonicity and show why its complete-data objective is
invariant under a change of common dominating measure.}
Use the setup and notation of \cref{ex:l5-continuous-em}. Fix
\(y_{\mathrm{obs}}\) such that, for every parameter under consideration,
\(m_\theta(y_{\mathrm{obs}})>0\) and
\(q_\theta(x,y_{\mathrm{obs}})>0\) for \(P_X\)-almost every \(x\). Put
\(r_k(x)=r_{\theta_k}(x\mid y_{\mathrm{obs}})\). Continuous EM uses
\[
 \theta_{k+1}\in\arg\max_\theta
 \int_\cX r_k(x)\log q_\theta(x,y_{\mathrm{obs}})\,P_X(dx).
\]
\begin{enumerate}[(a),beginpenalty=10000]
\item Prove the density analogue of the identity in \cref{ex:l5-em} and
deduce that this update does not decrease
\(\log m_\theta(y_{\mathrm{obs}})\).
\item Explain why the reference measure must be common to all parameter
values. More precisely, suppose another common reference measure satisfies
\(dP'=h\,dP\), where \(h>0\). Show that replacing \(P\) by \(P'\)
changes the expected complete-data log density by a term independent of
\(\theta\). Explain how counting measures recover the finite-space formulas
in \cref{ex:l5-em}.
\end{enumerate}
\end{exercisechunk}
